% ============================================================================
% Section 3 -- the capacity model. Contribution 2.
%
% COMPRESSED 2026-09-27. A 4-page analytical section in a 10-page systems paper
% pushed the evaluation to p11; the norm puts it near p7. The derivations,
% fits, residual structure and sensitivity now live in appendix A. What stays
% is what a reader needs to USE the model: the inputs it takes, the two
% conditions it applies, and what it returns.
% ============================================================================

\section{When does a second tier pay?}
\label{sec:model}

The model takes a profile of the device and returns the host cohort that maximises
decode throughput --- or zero, when this configuration should not use a second tier
at all. Below are the profile it needs, the two conditions it applies, and what its
structure says about which configurations are worth serving this way.

\subsection{Setup}

Let $\Xlink$ be the achieved interconnect bandwidth, $\HBM$ the device memory, $W$
the weight bytes, and $\Ckv = \HBM - W - \mathcal{R}$ the memory left for KV after a
$4$\,GB device reserve for activations and fragmentation. Let $S$ be the context
length and $\btok$ the KV bytes per token.

A decode step reads every weight once and every resident KV byte once, and the two
reads do not achieve the same bandwidth. Writing $\BW_w$ and $\BW_{kv}$ for each,
%
\begin{equation}
\label{eq:step}
\mathrm{ITL}(B) \;=\; \frac{W}{\BW_w} \;+\; \frac{B \cdot S \cdot \btok}{\BW_{kv}}.
\end{equation}
%
That weights stream faster than paged KV is unsurprising; we separate them
because the gap is large enough to change the answer (\S\ref{sec:model:params}).

The capacity-limited batch is $B_{\mathrm{cap}} = \Ckv/(S\btok)$. A step of length
$\mathrm{ITL}$ cannot carry more than $\mathrm{ITL}\cdot\Xlink$ bytes across the
link, so the host cohort is bounded by
%
\begin{equation}
\label{eq:boff}
B_{\mathrm{off}} \;=\; \frac{\mathrm{ITL} \cdot \Xlink}{S \btok}.
\end{equation}
%
That is a capacity bound and nothing more: it says how many riders' bytes fit, not
what they cost. What they cost is two numbers we measure --- a \textbf{fixed} $F$
for using the host path at all, whatever the rider count, and a
\textbf{per-request} $c$, the part of a rider's transfer that does not hide behind
the compute it is prefetched against. On A100/PCIe at 2048 context, $F = 12.3$\,ms
and $c = 2.2$--$2.9$\,ms (\S\ref{sec:eval:cost}).

\subsection{Two consequences that run against intuition}
\label{sec:model:cancel}

\paragraph{Context length cancels.} The ratio of riders to residents,
%
\begin{equation}
\label{eq:cancel}
\frac{B_{\mathrm{off}}}{B_{\mathrm{cap}}}
  \;=\; \frac{\mathrm{ITL}\cdot \Xlink}{\Ckv},
\end{equation}
%
is independent of $S$. Long context does not make offloading more attractive: what
it buys a rider in freed HBM it spends again in per-step transfer. This is worth
stating because the intuition runs the other way --- long context is what makes KV
expensive, so it looks like the regime where relocating it should pay most.

Equation~\ref{eq:step} claims ITL depends on batch and context only through their
product, and that the intercept is the weight term. On an A100-80 over 41
configurations both hold \emph{within} each model ($R^2 = 0.975$ and $0.92$), and
the intercept --- which the equation predicts rather than fits --- comes out exactly
for Qwen2.5-32B and 20\% high for Llama-3.1-8B; Appendix~\ref{app:params} gives the
fits and why the 8B misses. Pooling the models is meaningless because the intercept
is model-dependent, which is next.
% src: paper/model/p4_llama8b_a100.json, p4_qwen32b_a100.json @ 2bfc6ca5e

\paragraph{The step time is model-dependent.}
\label{sec:model:modeldep}
At capacity a step reads all of device memory once, which invites the
simplification $\mathrm{ITL}_0 \approx \HBM/\BW$ --- a quantity fixed by the device
alone. It does not hold. Evaluating Eq.~\ref{eq:step} at each model's own pool
gives $78.4$\,ms for Llama-3.1-8B against $53.2$\,ms for Qwen2.5-32B, a
$1.47\times$ spread on identical hardware, because a model that is mostly weights
reads device memory faster than one that is mostly KV --- so the gate below is
model-dependent rather than a property of the device. Both are
\emph{projections}, the largest pools fitting on an 80\,GB part measuring $69.7$
and $50.9$\,ms.
% src: paper/model/p4_llama8b_a100.json @ 2bfc6ca5e
% src: paper/model/p4_qwen32b_a100.json @ 2bfc6ca5e

\subsection{The gain law}
\label{sec:model:gain}

Adding $B_{\mathrm{off}}$ riders to $B_{\mathrm{cap}}$ residents multiplies the
served batch by
%
\begin{equation}
\label{eq:gain-exact}
\text{gain} \;=\; 1 + \frac{B_{\mathrm{off}}}{B_{\mathrm{cap}}}
      \;=\; 1 + \frac{\mathrm{ITL}_0 \cdot \Xlink}{\Ckv},
\end{equation}
%
exact in the fluid limit, and the form to compute with. Its structure is clearer on
noting that a full pool takes about $\HBM$ over the bandwidth the device achieves
across it, so the gain is roughly $(\Xlink/\BW_{\mathrm{pool}})\cdot(\HBM/\Ckv)$.
The bandwidth ratio is a handicap --- $419$\,GB/s against a pool reading at
$2000$--$2900$ on a GH200 --- that nothing a serving system improves. The amplifier
is $\HBM/\Ckv$: \textbf{the value of a second tier tracks how much of memory the
weights ate}, not model size and not context length --- the section's most useful
consequence, because it runs counter to how the field picks evaluation models.

Equation~\ref{eq:gain-exact} counts requests the device can \emph{hold}, not
requests per second it serves. Each rider also displaces a share of a resident's
blocks for the window it streams into, which is why Table~\ref{tab:gain}'s
capacity column ($\times1.50$ for Qwen2.5-32B fp16) sits below a hand
recomputation ($\times1.54$).

\paragraph{Weight precision is the clean lever.} Comparing models of different
sizes moves the weight footprint but also layer count, KV geometry and
architecture; the clean instrument is one model at two precisions, and the two
Qwen2.5-32B rows differ in nothing else. \textbf{Neither factor of the multiplier
is measured}: the bandwidth ratio needs two interconnect generations and we have
one, and the amplifier needs that precision pair run end to end. The rows are
projections.

\begin{table*}[t]
\centering
\small
\begin{tabular}{lrrrrrrr}
\toprule
Model (GH200-96) & weights & $\Ckv$ & $\HBM/\Ckv$ & ITL$_0$ & $F$ & capacity & tput \\
\midrule
Llama-8B fp16          & 16.1\,GB & 75.9\,GB & 1.26 & 46.6\,ms & \phantom{0}6.2\,ms & $\times$1.22 & $\times$0.98 \\
Qwen-32B \textbf{fp8}  & 32.0\,GB & 60.0\,GB & 1.60 & 43.3\,ms & 12.3\,ms & $\times$1.28 & $\times$0.91 \\
Qwen-32B \textbf{fp16} & 65.5\,GB & 26.5\,GB & 3.63 & 33.9\,ms & 12.3\,ms & $\times$1.50 & $\times$\textbf{1.01} \\
Llama-70B fp8          & 72.7\,GB & 19.3\,GB & 4.97 & 33.0\,ms & 15.4\,ms & $\times$1.67 & $\times$\textbf{1.05} \\
\bottomrule
\end{tabular}
\caption{\emph{Projections}, not measurements, at 2048 context and an 80\,ms target
--- the first round figure above $\mathrm{ITL}_0 + F$ for every row.
$\Ckv = \HBM - W - 4$\,GB. \emph{capacity} charges each rider the blocks it
displaces and the two HBM passes its KV makes; \emph{tput} additionally charges $F$
and the per-request cost. \textbf{The two rows where weights consume most of the
device pay; the two where they do not, do not} --- the amplifier doing the work.
\textbf{The two Qwen-32B rows differ only in weight precision.} Weights are
measured checkpoints except Qwen-32B fp8; Appendix~\ref{app:params} gives the
inputs and the sensitivity.}
% src: paper/model/tables.py (wraps killer_v2.py:analyze, S=8192, GH200-96 at
% the measured 419 GB/s, slo=80 ms, stage=2, F per layer, c = 12% of link)
\label{tab:gain}
\end{table*}

\subsection{From capacity to throughput: the second gate}
\label{sec:model:omits}

A serving system is paid in requests \emph{per second}. Charging the host path's
two costs gives
%
\begin{equation}
\label{eq:gain-tput}
\text{gain}_{\text{tput}} \;=\;
  \Big(1 + \tfrac{B_{\mathrm{off}}}{B_{\mathrm{cap}}}\Big)
  \cdot \frac{\mathrm{ITL}_0}{\mathrm{ITL}_0 + F + \Delta},
\end{equation}
%
where $\Delta = c\,B_{\mathrm{off}}$. This is a \textbf{second gate} and it binds
harder than the latency one: the technique pays only when
$B_{\mathrm{off}}/B_{\mathrm{cap}} > (F + \Delta)/\mathrm{ITL}_0$.

\textbf{Two of Table~\ref{tab:gain}'s four rows clear it.} Charging each the $F$ it
would pay ($12.3$\,ms measured on 64 layers, charged per layer, since
\S\ref{sec:impl:cost} attributes it to per-layer launch structure), capacity gains of
$\times$1.22 to $\times$1.67 become throughput ratios of $\times$0.98, $\times$0.91,
$\times$1.01 and $\times$1.05 --- the ordering the amplifier predicts. The margins
are thin, so the useful way to report them is to invert Eq.~\ref{eq:gain-tput} for
the fixed cost each configuration can afford:
%
\begin{equation}
\label{eq:fstar}
F^{*} \;=\; \Big(\tfrac{B_{\mathrm{tot}}}{B_0} - 1\Big)\mathrm{ITL}_0 - \Delta.
\end{equation}
%
Eq.~\ref{eq:fstar} gives $F^{*} = 12.8$\,ms for Qwen2.5-32B fp16 against the
$12.3$ it pays, and $18.2$ against $15.4$ for Llama-70B fp8 --- margins of $0.5$
and $2.8$\,ms; Llama-3.1-8B and Qwen2.5-32B fp8 miss, needing $23\%$ and $44\%$
less $F$. \textbf{So the paper's positive claim is a margin of a few milliseconds
in one measured quantity}, which is why measuring $F$ on a coherent link is the
first experiment we ask of one: it is launch-bound rather than bandwidth-bound, so
a faster link does not touch it, but it scales with the host CPU and layer depth.

Four smaller terms are omitted --- host-side writes, host DRAM bandwidth, the
integrality of both cohorts, and MoE weight reads --- and \S\ref{sec:discussion}
states which regimes each of them bounds.

\subsection{The applicability gate}
\label{sec:model:gate}

Capacity binds only when the latency target exceeds the full-pool step time
\emph{plus} the fixed cost of using the host path:
%
\begin{equation}
\label{eq:gate}
\mathrm{SLO} \;>\; \frac{W}{\BW_w} + \frac{\Ckv}{\BW_{kv}} + F.
\end{equation}
%
Below it the scheduler is already refusing requests HBM could hold, because
admitting them would break latency: the system is latency-limited rather than
capacity-limited, and offloading is strictly harmful. Its effect is continuous
rather than binary --- a target only slightly above $\mathrm{ITL}_0 + F$ leaves room
for very few riders whatever the link delivers, which is why Table~\ref{tab:gain} is
computed above that sum for every row. It scopes the technique to
\emph{throughput-oriented} serving and excludes interactive chat, and also excludes
latency-sensitive agentic pipelines, which are often mis-filed as throughput work
because their requests are long even when a user is waiting.

\paragraph{Measured parameters.}
\label{sec:model:params}
Both bandwidths are measured on A100-80 across 41 configurations on two models;
$\BW_w$ is stable within a model at about $0.75$ of the measured $1935$\,GB/s peak.
$\BW_{kv}$ is \textbf{not} constant: it degrades by $1.6\times$ as the KV working
set grows, and worst exactly at the capacity-bound operating point this paper is
about --- so \textbf{a roofline charging KV at one fixed bandwidth is optimistic
precisely where it matters}. We measure that curve on Llama-3.1-8B alone, and
Table~\ref{tab:gain} applies it to working sets of 19--76\,GB; the curve, the fits
and the sensitivity are in Appendix~\ref{app:params}.

\paragraph{A published result the model is consistent with.}
\label{sec:model:directkv}
DirectKV~\cite{directkv} reports its host-KV attention slightly slower than an
HBM-resident baseline on Llama-3.1-8B and winning on OPT-30B, whose 60\,GB of
weights triple the amplifier. We put it no more strongly than \emph{consistent
with}: Appendix~\ref{app:directkv} gives three objections, including that OPT-13B's
amplifier is close to the 8B's yet they win on it.
